\chapter{Estados, reversibilidad y cinemática}

La mecánica clásica no es solo un conjunto de ecuaciones para cuerpos que caen,
planetas en órbita o partículas cargadas. Antes de elegir una ecuación concreta,
podemos plantear una pregunta más general: ¿qué aspecto debe tener una ley del
movimiento para considerarse clásica? Empezaremos con sistemas tan sencillos que
podemos dibujar en la pizarra todas sus historias posibles. Una vez aclarados allí
los conceptos de estado, determinismo, reversibilidad y conservación, volveremos a
las partículas que se mueven en el espacio ordinario.

\section{Dos tipos de preguntas}

En mecánica hay dos tipos de preguntas. Una busca la ley específica que rige un
sistema particular. Un planeta que se mueve en el campo gravitatorio de un cuerpo
masivo obedece una ley; una partícula con carga eléctrica que se mueve en un
campo magnético obedece otra. Son preguntas sobre las leyes detalladas de la
naturaleza.

La otra pregunta se refiere al marco común a todos esos ejemplos. ¿Hay reglas
que deba cumplir toda ley admisible? ¿Qué características pertenecen a la propia
mecánica clásica, con independencia de la fuerza o del objeto estudiado? Los
ejemplos concretos ilustrarán la respuesta, pero es esta segunda pregunta la que
organiza la materia.

El primer ejemplo más útil es una moneda abstracta sobre una mesa abstracta.
Conservamos una sola propiedad de la moneda: puede mostrar cara o cruz. Por tanto,
su espacio de configuraciones es
\begin{equation}
  \mathcal C=\{H,T\}.
\end{equation}
La palabra \emph{configuración} designa una de las formas posibles en que
puede encontrarse el sistema. Una condición inicial consiste en elegir una de
essas configuraciones. Una moneda con una sola cara sería aún más sencilla, pero
solo tendría una configuración y nunca podría ocurrir nada. Dos estados bastan
para que haya dinámica.

\section{La moneda: estado, ley e historia}

Hagamos avanzar el tiempo en pasos enteros, como si el mundo estuviera iluminado
por un estroboscopio:
\begin{equation}
  t=\ldots,-2,-1,0,1,2,\ldots .
\end{equation}
Una ley del movimiento es una regla de actualización. Asigna al estado en el
instante \(t\) un estado en el instante \(t+1\). La regla más sencilla no hace
nada:
\begin{equation}
  H\longmapsto H,
  \qquad
  T\longmapsto T.
  \label{eq:coin-identity}
\end{equation}
Si partimos de \(H\), la historia completa es \(H,H,H,\ldots\); si partimos de
\(T\), es \(T,T,T,\ldots\). La regla es aburrida, pero permite hacer predicciones
completas.

La otra regla reversible intercambia los dos estados:
\begin{equation}
  H\longmapsto T,
  \qquad
  T\longmapsto H.
  \label{eq:coin-flip}
\end{equation}
Ahora las historias alternan:
\begin{equation}
  H,T,H,T,\ldots,
  \qquad
  T,H,T,H,\ldots .
\end{equation}
Una flecha apunta del estado anterior a la actualización al estado posterior.
Así, el origen de la flecha indica dónde empieza el sistema y la punta, adónde
va.

\begin{figure}[htbp]
  \centering
  \includegraphics[width=0.86\textwidth]{lecture_01_figure_02.png}
  \caption{Las dos configuraciones de la moneda y la ley de inversión dirigida
  en la pizarra de la clase. La regla de actualización limpia aparece en
  la ec.~\eqref{eq:coin-flip}.}
  \lectureframe{00:06:50}
  \label{fig:coin-board}
\end{figure}

Podemos codificar las caras mediante una variable de dos valores:
\begin{equation}
  \sigma=
  \begin{cases}
    +1,&H,\\
    -1,&T.
  \end{cases}
\end{equation}
En física se suele utilizar la letra \(\sigma\) para magnitudes de dos valores,
especialmente al describir el espín, aunque aquí no hace falta saber nada sobre
el espín. Con esta notación, las dos reglas se expresan como ecuaciones de
movimiento:
\begin{align}
  \sigma(t+1)&=\sigma(t),
  &&\text{identity law},\label{eq:sigma-identity}\\
  \sigma(t+1)&=-\sigma(t),
  &&\text{flip law}.\label{eq:sigma-flip}
\end{align}
La variable no añade física nueva. Proporciona una etiqueta compacta para la
configuración y convierte un grafo en una ecuación.

\begin{figure}[htbp]
  \centering
  \includegraphics[width=0.86\textwidth]{lecture_01_figure_03.png}
  \caption{La ley de inversión escrita como
  \(\sigma(t+1)=-\sigma(t)\). Los valores enmarcados \(\sigma=\pm1\) identifican cara y
  cruz.}
  \lectureframe{00:08:58}
  \label{fig:sigma-board}
\end{figure}

Para cualquiera de las dos reglas, el estado inicial y la ley determinan el
estado en todo instante posterior. Esta afirmación presupone que el sistema está
\emph{cerrado}. Si alguien mete la mano y le da la vuelta a la moneda,
entonces esa mano es una influencia externa que no está incluida en nuestra
descripción. No hemos refutado el determinismo; hemos cambiado el sistema. Una
ley clásica pretende regir un sistema cerrado, con todas las influencias
pertinentes para su evolución incluidas.

Hay cuatro funciones matemáticas de un conjunto de dos elementos en sí mismo.
Las dos que no hemos escrito arriba envían ambas entradas a \(H\) o ambas a
\(T\). Son deterministas hacia el futuro, pero borran la distinción entre dos
pasados posibles. Cuando más adelante impongamos la reversibilidad, las
ec.~\eqref{eq:coin-identity} y \eqref{eq:coin-flip} serán exactamente las dos aplicaciones
permitidas.

\section{Un dado: órbitas y cantidades conservadas}

Sustituyamos la moneda por un dado. Para lo que nos interesa aquí, el único
hecho importante sobre un dado es que tiene seis configuraciones:
\begin{equation}
  \mathcal C=\{1,2,3,4,5,6\}.
\end{equation}
Una condición inicial consiste en elegir uno de los números. Podríamos hacer que
todos los estados permanecieran fijos, pero los grafos dirigidos son más útiles
para reglas menos triviales. Consideremos
\begin{equation}
  1\longmapsto2\longmapsto3\longmapsto4
  \longmapsto5\longmapsto6\longmapsto1.
  \label{eq:six-cycle}
\end{equation}
Si empezamos en \(4\), por ejemplo, la historia es

\[
  4,5,6,1,2,3,4,\ldots .
\]
Los seis estados forman un ciclo.\footnote{Nota de la fuente: durante la primera
lectura de las flechas, una frase pronunciada dice que \(6\) vuelve a \(6\). La
frase inmediatamente posterior, el ciclo dibujado y la historia enumerada dan
todos la actualización pretendida \(6\mapsto1\).}

\begin{figure}[htbp]
  \centering
  \includegraphics[width=0.86\textwidth]{lecture_01_figure_06.png}
  \caption{El único ciclo de seis estados. En cada paso estroboscópico, un
  estado pasa al siguiente y vuelve al cabo de seis actualizaciones.}
  \lectureframe{00:12:40}
  \label{fig:six-cycle-board}
\end{figure}

Una regla puede parecer más complicada sin ser estructuralmente distinta. Por
ejemplo,
\begin{equation}
  1\to2\to5\to3\to4\to6\to1
  \label{eq:relabeled-cycle}
\end{equation}
sigue siendo un único ciclo de seis estados. Si colocamos los vértices en la
página en el orden \(1,2,5,3,4,6\), obtenemos el mismo dibujo que en la
ec.~\eqref{eq:six-cycle}. Etiquetas como \(1\) y \(5\) son nombres; la estructura del
ciclo no depende de ellas.

Una ley realmente distinta divide los seis estados en dos ciclos, por ejemplo:
\begin{equation}
  1\to2\to3\to1,
  \qquad
  4\to6\to5\to4.
  \label{eq:two-cycles}
\end{equation}
Sigue habiendo un solo dado y una sola ley, pero una trayectoria que empieza en
un ciclo nunca puede llegar al otro. Podemos asignar a los ciclos un número
\(Q\):
\begin{equation}
  Q=
  \begin{cases}
    0,&\text{on the cycle }\{1,2,3\},\\
    1,&\text{on the cycle }\{4,5,6\}.
  \end{cases}
\end{equation}
El estado cambia en cada paso, mientras que \(Q\) no cambia.

\begin{definition}
Una cantidad conservada es una función del estado cuyo valor no cambia bajo la
ley del movimiento.
\end{definition}

La misma idea admite un punto fijo, un ciclo de dos estados y otro de tres. La
etiqueta conservada indica qué órbita ocupa el sistema; no tiene por qué
identificar el estado concreto dentro de esa órbita. Así, la conservación es
compatible con el movimiento. Lo que se conserva no es necesariamente la
configuración completa, sino alguna distinción entre las configuraciones
posibles.

\begin{figure}[htbp]
  \centering
  \includegraphics[width=0.86\textwidth]{lecture_01_figure_07.png}
  \caption{Varios ciclos desconectados en la pizarra: un punto fijo, un ciclo
  de dos estados y otro de tres. La etiqueta del ciclo se conserva, aunque el
  estado recorra ese ciclo.}
  \lectureframe{00:18:40}
  \label{fig:conserved-cycles-board}
\end{figure}

En un sistema reversible finito, toda trayectoria es un ciclo. Un único ciclo
no admite una etiqueta conservada no constante de este tipo elemental, porque la
trayectoria acaba visitando todos los estados. Varios ciclos sí permiten tal
etiqueta, ya que la evolución nunca lleva un estado de una componente a otra.

\section{¿Qué debe incluirse en un estado?}

Los ejemplos de la moneda usaban una regla de un paso: la cara siguiente dependía
solo de la cara actual. Una pregunta del público pone a prueba si esa restricción
es esencial.

\begin{classroomqa}
\audiencequestion{¿Se podría construir una secuencia de caras y cruces en la que
el resultado siguiente dependiera de algo más que de la cara visible en ese
momento, quizá de uno o varios resultados anteriores?\lecturetimestamp{00:19:27--00:21:23}}

\lecturerresponse{Sí. En ese caso, \(H\) o \(T\) por sí solo no es el estado
completo. Si hace falta conocer la cara anterior, ampliemos el espacio de estados
a
\(
  \{HH,HT,TH,TT\}.
\)
Cada par ordenado registra la cara actual y la cara anterior. Así, la evolución
puede expresarse de nuevo como una regla de un paso en este espacio de cuatro
estados. Si hace falta aún más historia, incluyamos también esa información.}
\end{classroomqa}

Esta respuesta da el significado operativo de estado.

\begin{definition}
El estado de un sistema en un instante consta de toda la información necesaria,
junto con la ley del movimiento, para predecir su futuro.
\end{definition}

Por tanto, el estado no es simplemente una fotografía. Es el conjunto completo
de datos que hace definida la evolución. Más adelante, al describir una partícula
newtoniana, su posición no bastará; los datos iniciales incluirán también la
velocidad. La moneda con memoria es la versión más sencilla del mismo principio.

\section{Espacios de estados infinitos}

Un número infinito de estados no requiere un número infinito de objetos.
Consideremos una partícula confinada a los sitios enteros de una recta infinita.
Su estado es un entero:
\[
  n\in\mathbb Z.
\]
Una ley la desplaza un sitio a la derecha en cada paso:
\begin{equation}
  n(t+1)=n(t)+1.
  \label{eq:lattice-one}
\end{equation}
Esta órbita nunca se cierra, pero, si se sigue tanto hacia el futuro como hacia
el pasado, pasa por todos los enteros. En ese sentido extendido, desempeña el
papel de un ciclo infinito.

Otra ley la desplaza dos sitios cada vez:
\begin{equation}
  n(t+1)=n(t)+2.
  \label{eq:lattice-two}
\end{equation}
Ahora, los sitios pares evolucionan solo entre sitios pares, y los impares, solo
entre sitios impares. La paridad se conserva:
\begin{equation}
  \begin{aligned}
    Q_{\rm parity}(n)&=n\bmod 2,\\
    Q_{\rm parity}(n(t+1))&=Q_{\rm parity}(n(t)).
  \end{aligned}
\end{equation}
El paso a un espacio de configuraciones infinito permite historias que nunca se
repiten, pero no cambia el significado de estado, órbita o conservación.

\section{El determinismo aún no es reversibilidad}

Ahora dibujaremos una regla que predice de forma única el futuro, pero que no
cumple un segundo requisito clásico. Demos a una moneda de tres caras los estados
cara, cruz y canto, y definamos
\begin{equation}
  H\to T,
  \qquad
  T\to E,
  \qquad
  E\to T.
  \label{eq:nonreversible-map}
\end{equation}
Si empezamos en \(H\), la historia es
\[
  H,T,E,T,E,\ldots;
\]
Si empezamos en \(T\) o en \(E\), la alternancia posterior también está
determinada sin ambigüedad. Desde cada estado sale exactamente una flecha, así que
la regla es determinista hacia el futuro.

La dificultad aparece cuando intentamos reconstruir el pasado. Si el estado
actual es \(T\), el estado anterior podría haber sido \(H\) o \(E\). Dos
historias han confluido. Si invertimos todas las flechas, el defecto resulta
inmediato: desde \(T\), el grafo invertido ofrece dos movimientos posibles. La
regla inversa no es determinista.

\begin{figure}[htbp]
  \centering
  \includegraphics[width=0.86\textwidth]{lecture_01_figure_08.png}
  \caption{La ley no reversible \(H,T,E\) y algunas historias representativas.
  Tanto \(H\) como \(E\) conducen a \(T\), por lo que el presente no
  determina un pasado único.}
  \lectureframe{00:28:50}
  \label{fig:nonreversible-board}
\end{figure}

\begin{definition}
Una ley discreta es reversible cuando el estado presente determina tanto un
estado siguiente único como un estado anterior único. En el grafo de estados,
de cada estado debe salir exactamente una flecha y a cada estado debe llegar
exactamente una flecha.
\end{definition}

En un espacio de estados finito, una actualización reversible es una permutación
de los estados. No crea ni destruye distinciones entre estados; solo los
reordena. Por eso su grafo se descompone en ciclos. La aplicación de la
ec.~\eqref{eq:nonreversible-map} es una función, pero no es inyectiva y, por tanto, no tiene
inversa.

\begin{classroomqa}
\audiencequestion{¿Se podría admitir una regla de este tipo si los estados
representaran probabilidades de una partícula?\lecturetimestamp{00:29:46--00:31:17}}

\lecturerresponse{Eso correspondería a un marco distinto. La mecánica
clásica fundamental, tal como la hemos formulado aquí, utiliza leyes
deterministas en vez de reglas de actualización probabilísticas. Sin duda se
puede inventar y estudiar una regla así, pero en este sentido no es una ley de la
mecánica clásica. La mecánica cuántica no es determinista del mismo modo y, aun
así, tiene un análogo de la evolución reversible; para tratarlo hace falta una
discusión posterior. La mecánica estadística clásica también utiliza
probabilidades, pero parte de leyes mecánicas subyacentes y representa el
conocimiento incompleto de los datos iniciales.}
\end{classroomqa}

La distinción es importante. Una distribución de probabilidad que se usa porque
desconocemos el estado clásico exacto no es una ley del movimiento indeterminista.
Describe lo que sabemos de un sistema cuya evolución microscópica sigue
considerándose determinista y reversible.

\section{Leyes exactas y conocimiento imperfecto}

La predicción requiere dos elementos:
\begin{enumerate}
  \item la ley del movimiento;
  \item el estado inicial.
\end{enumerate}
Una ley exacta no produce una predicción exacta si la condición inicial solo se
conoce de forma aproximada. En el caso de la moneda, es fácil imaginar que
miramos la mesa y vemos perfectamente qué cara ha salido. Las variables
mecánicas reales son continuas. Las posiciones y velocidades pueden tomar
cualquier valor real, y ninguna medición proporciona infinitas cifras decimales.

Supongamos que el estado inicial verdadero se encuentra dentro de una región
pequeña de estados posibles. Con algunas dinámicas, esa región permanece pequeña
durante mucho tiempo; con otras, los estados iniciales próximos se separan
rápidamente. En un sistema caótico, las pequeñas incertidumbres iniciales acaban
produciendo grandes diferencias. La atmósfera es un ejemplo conocido. Mejores
mediciones pueden prolongar el intervalo durante el cual la predicción resulta
útil, pero para cualquier precisión de medida fija existe un instante lo bastante
lejano en que la predicción deja de ser fiable.

Esta limitación práctica no introduce ambigüedad en las ecuaciones. Si las
condiciones iniciales se conocieran con exactitud, una ley clásica reversible
determinaría tanto el futuro como el pasado. Con precisión finita, la pregunta
cuantitativa pasa a ser: ¿con qué precisión hay que conocer el estado inicial
para predecir durante un intervalo de tiempo especificado?

La prueba gráfica nos lleva ahora a una conclusión concisa. Una ley representada
por una flecha entrante y otra saliente en cada estado es determinista y
reversible. Las afirmaciones correspondientes para las partículas serán menos
elementales, pero expresan la misma estructura clásica.

\section{De los estados abstractos a los puntos del espacio}

Pasemos ahora de las monedas y los dados a las partículas puntuales. Una
partícula puede moverse en una, tres, cinco dimensiones o en cualquier otro
número de dimensiones, pero el espacio ordinario nos da un ejemplo inmediato.
Para describir cuantitativamente un punto, elijamos:
\begin{itemize}
  \item un origen;
  \item tres ejes mutuamente perpendiculares;
  \item una orientación para esos ejes;
  \item una unidad de longitud.
\end{itemize}
La elección del origen es convencional. Se puede situar en Palo Alto o en
cualquier otro lugar, y las conclusiones físicas no deberían depender de esa
elección. Los ejes se pueden etiquetar \(x,y,z\) o \(x_1,x_2,x_3\). Una vez fijadas las
direcciones \(x\) y \(y\), el eje \(z\) puede apuntar en cualquiera de
dos direcciones perpendiculares. Adoptamos la convención dextrógira: con el pulgar
de la mano derecha apuntando en la dirección \(x\) y el índice en la
dirección \(y\), el dedo medio apunta en la dirección \(z\).

Para marcar intervalos iguales en los ejes hace falta una unidad, como metros,
pies, pulgadas o años luz. Una vez fijadas estas convenciones, cada punto queda
identificado por tres números:
\begin{equation}
  (x,y,z).
\end{equation}
Se puede llegar al punto desde el origen mediante desplazamientos a lo largo de
los tres ejes. Un sistema de muchas partículas requiere una terna de este tipo
para cada partícula.

\begin{figure}[htbp]
  \centering
  \includegraphics[width=0.86\textwidth]{lecture_01_figure_09.png}
  \caption{Ejes cartesianos y las tres coordenadas de un punto. El dibujo
  fija una convención de coordenadas; el punto en sí no depende del lugar donde
  decidamos situar el origen o los ejes.}
  \lectureframe{00:43:00}
  \label{fig:coordinates-board}
\end{figure}

\section{Vectores y sus componentes}

Un vector es un objeto con módulo y dirección. El desplazamiento desde el
origen hasta un punto es un vector, pero los vectores también describen la
velocidad, la aceleración, los campos eléctricos y muchas otras magnitudes.
Escribiremos los vectores en negrita, por ejemplo \(\mathbf v\), y su módulo como
\(\lvert\mathbf v\rvert\). El módulo es no negativo y solo se anula para el
vector cero.

Geométricamente, un vector es una flecha libre. Trasladar la flecha sin cambiar
su longitud ni su dirección no altera el vector. En una aplicación concreta, un
vector puede estar asociado a un punto, como cuando especificamos el campo
eléctrico en una posición determinada, pero esa asociación es información
adicional y no forma parte del vector abstracto en sí.

Con respecto a los ejes cartesianos, un vector tiene componentes:
\begin{equation}
  \mathbf v=(v_x,v_y,v_z).
\end{equation}
Las componentes son números ordinarios con signo, que se obtienen al proyectar el
vector sobre los ejes. Codifican tanto el módulo como la dirección. El teorema de
Pitágoras en tres dimensiones da
\begin{equation}
  \lvert\mathbf v\rvert^2
  =v_x^2+v_y^2+v_z^2,
  \qquad
  \lvert\mathbf v\rvert
  =\sqrt{v_x^2+v_y^2+v_z^2}.
  \label{eq:vector-magnitude}
\end{equation}
Las razones entre las componentes determinan la dirección. Una componente
negativa significa que el vector apunta parcialmente en el sentido negativo de
ese eje.

\subsection{Multiplicación por un escalar y suma}

Multiplicar un vector por un número real cambia su longitud y, si el número es
negativo, invierte su dirección. Para \(c>0\), \(c\mathbf A\) apunta en la dirección
de \(\mathbf A\) y su módulo es
\(c\lvert\mathbf A\rvert\). En particular, \(-\mathbf A\) tiene la misma
longitud que \(\mathbf A\) y apunta en sentido contrario.

Hay tres representaciones equivalentes de la suma de vectores. Primero,
coloquemos el origen de \(\mathbf B\) en la punta de \(\mathbf A\); la flecha que va
desde el origen de \(\mathbf A\) hasta la punta de \(\mathbf B\) es
\(\mathbf A+\mathbf B\).
Segundo, coloquemos juntos los orígenes de ambos vectores y completemos un
paralelogramo; su diagonal es la suma. Esta representación hace patente la
conmutatividad:
\begin{equation}
  \mathbf A+\mathbf B=\mathbf B+\mathbf A.
\end{equation}
Tercero, sumemos las componentes correspondientes:
\begin{equation}
  C_i=A_i+B_i,
  \qquad i=x,y,z,
  \label{eq:component-addition}
\end{equation}
donde \(\mathbf C=\mathbf A+\mathbf B\). Estas reglas también definen combinaciones
lineales arbitrarias, como \(a\mathbf A+b\mathbf B\).

\begin{figure}[htbp]
  \centering
  \includegraphics[width=0.86\textwidth]{lecture_01_figure_10.png}
  \caption{Construcciones del triángulo y del paralelogramo para la suma de
  vectores. Ambas dan el mismo vector \(\mathbf A+\mathbf B\).}
  \lectureframe{00:53:00}
  \label{fig:vector-addition-board}
\end{figure}

En general, no existe una operación para dividir un vector por otro. Sí hay,
sin embargo, dos productos vectoriales estándar. Dejaremos el producto vectorial
para una clase posterior. Aquí necesitaremos el producto escalar.

\subsection{El producto escalar}

Proyectemos \(\mathbf A\) sobre el eje que define \(\mathbf B\). Si el ángulo entre los
vectores es \(\theta\), la componente con signo de \(\mathbf A\) en la dirección de
\(\mathbf B\) es \(\lvert\mathbf A\rvert\cos\theta\). Al multiplicar esta
proyección por \(\lvert\mathbf B\rvert\) se define
\begin{equation}
  \mathbf A\cdot\mathbf B
  =\lvert\mathbf A\rvert\,\lvert\mathbf B\rvert\cos\theta.
  \label{eq:dot-geometric}
\end{equation}
El resultado es un escalar, no otro vector. La ec.~\eqref{eq:dot-geometric}
es simétrica con respecto a \(\mathbf A\) y \(\mathbf B\), de modo que
\(\mathbf A\cdot\mathbf B=\mathbf B\cdot\mathbf A\).

\begin{figure}[htbp]
  \centering
  \includegraphics[width=0.86\textwidth]{lecture_01_figure_11.png}
  \caption{Construcción del producto escalar mediante proyección en la pizarra.
  La expresión tipográfica es
  \(\mathbf A\cdot\mathbf B=AB\cos\theta\).}
  \lectureframe{00:58:30}
  \label{fig:dot-product-board}
\end{figure}

El signo contiene información geométrica. El producto escalar es positivo si
el ángulo es agudo, negativo si es obtuso y cero si los vectores son
perpendiculares. Por tanto,
\begin{equation}
  \mathbf A\perp\mathbf B
  \quad\Longleftrightarrow\quad
  \mathbf A\cdot\mathbf B=0
\end{equation}
para vectores no nulos.

En componentes,
\begin{equation}
  \mathbf A\cdot\mathbf B
  =A_xB_x+A_yB_y+A_zB_z.
  \label{eq:dot-components}
\end{equation}
Esta fórmula permite calcular directamente. En particular,
\begin{equation}
  \mathbf A\cdot\mathbf A
  =A_x^2+A_y^2+A_z^2
  =\lvert\mathbf A\rvert^2.
  \label{eq:self-dot}
\end{equation}
Conocidas las componentes de dos vectores, las ec.~\eqref{eq:dot-geometric} y
\eqref{eq:dot-components} determinan el ángulo entre ellos:
\begin{equation}
  \cos\theta
  =\frac{\mathbf A\cdot\mathbf B}
  {\lvert\mathbf A\rvert\,\lvert\mathbf B\rvert}.
\end{equation}

De la misma álgebra se deduce la ley de los cosenos. Sean \(\mathbf A\) y
\(\mathbf B\) los vectores que van de un vértice común a los otros dos vértices de un
triángulo. El tercer lado es
\(\mathbf C=\mathbf A-\mathbf B\), no \(\mathbf A+\mathbf B\). Por tanto,
\begin{align}
  C^2
  &=(\mathbf A-\mathbf B)\cdot(\mathbf A-\mathbf B)\notag\\
  &=A^2+B^2-2\mathbf A\cdot\mathbf B\notag\\
  &=A^2+B^2-2AB\cos\theta.
  \label{eq:law-of-cosines}
\end{align}

\begin{figure}[htbp]
  \centering
  \includegraphics[width=0.86\textwidth]{lecture_01_figure_04.png}
  \caption{El triángulo utilizado para la ley de los cosenos. La corrección
  en la pizarra identifica el lado que cierra el triángulo como
  \(\mathbf C=\mathbf A-\mathbf B\), lo que conduce a la
ec.~\eqref{eq:law-of-cosines}.}
  \lectureframe{01:04:29}
  \label{fig:law-cosines-board}
\end{figure}

\section{Posición, velocidad y aceleración}

Elijamos un origen y sea \(\mathbf r\) el vector de posición desde el origen hasta
una partícula. Sus componentes son las coordenadas cartesianas:
\begin{equation}
  \mathbf r=(x,y,z),
  \qquad
  \lvert\mathbf r\rvert=\sqrt{x^2+y^2+z^2}.
\end{equation}
Cambiar el origen modifica \(\mathbf r\); pospondremos esa cuestión. Una vez fijado
el origen, el movimiento se describe mediante una posición dependiente del
tiempo:
\begin{equation}
  \mathbf r(t)=\bigl(x(t),y(t),z(t)\bigr).
\end{equation}
Así pues, la trayectoria queda codificada por tres funciones del tiempo.

La velocidad es la derivada temporal de la posición:
\begin{equation}
  \mathbf v(t)=\frac{d\mathbf r}{dt}
  =\dot{\mathbf r}(t),
  \qquad
  v_i=\frac{dx_i}{dt}=\dot x_i.
  \label{eq:velocity}
\end{equation}
Un punto siempre denota una derivada con respecto al tiempo. El módulo
\(\lvert\mathbf v\rvert\) es la rapidez. En general, la posición y la velocidad
son independientes en cuanto a sus direcciones: una partícula situada en un
lugar puede moverse en cualquier dirección permitida por el problema.

La aceleración es la tasa de cambio de la velocidad y, por tanto, la segunda
derivada temporal de la posición:
\begin{equation}
  \mathbf a(t)=\frac{d\mathbf v}{dt}
  =\dot{\mathbf v}(t)
  =\frac{d^2\mathbf r}{dt^2}
  =\ddot{\mathbf r}(t),
  \qquad
  a_i=\ddot x_i.
  \label{eq:acceleration}
\end{equation}
La aceleración también es un vector. Si una partícula que se mueve en la
dirección positiva de \(x\) aumenta su rapidez, \(a_x>0\). Si la reduce,
\(a_x<0\). Lo que suele llamarse desaceleración es simplemente una aceleración
opuesta a la velocidad actual. En general, la velocidad puede cambiar de
dirección aunque su módulo permanezca constante.

\begin{figure}[htbp]
  \centering
  \includegraphics[width=0.86\textwidth]{lecture_01_figure_05.png}
  \caption{Componentes de la velocidad expresadas como derivadas temporales de
  las componentes de la posición. En notación vectorial compacta,
  \(\mathbf v=d\mathbf r/dt\).}
  \lectureframe{01:14:24}
  \label{fig:velocity-board}
\end{figure}

\subsection{Aceleración uniforme en una recta}

Para el movimiento unidimensional, escribamos
\begin{equation}
  x(t)=a+bt+ct^2.
  \label{eq:quadratic-motion}
\end{equation}
En el instante \(t=0\), la partícula se encuentra en \(x=a\). Derivando una y
luego dos veces, obtenemos
\begin{align}
  v(t)=\dot x(t)&=b+2ct,\label{eq:line-velocity}\\
  a(t)=\ddot x(t)&=2c.\label{eq:line-acceleration}
\end{align}
Así, \(b\) es la velocidad inicial y \(2c\) es una aceleración constante. Si
\(c>0\), la velocidad aumenta linealmente; si \(c<0\), disminuye linealmente.
Esta ley cuadrática para la posición es la forma que adopta el movimiento bajo
una aceleración constante, como en el movimiento idealizado de un cuerpo que cae
cerca de la superficie terrestre.

\section{Movimiento circular uniforme}

Para el último ejemplo, hagamos que una partícula se mueva en el plano \(x\)-\(y\)
sobre la circunferencia unidad. Parte del eje \(x\) positivo en el instante
\(t=0\), y su coordenada angular aumenta uniformemente:
\begin{equation}
  \theta(t)=\omega t.
\end{equation}
La constante \(\omega\) es la frecuencia angular. Una revolución corresponde a un
aumento de \(2\pi\) en el ángulo, por lo que el período es
\begin{equation}
  T=\frac{2\pi}{\omega}.
  \label{eq:circular-period}
\end{equation}
Podemos permitir que el ángulo aumente más allá de \(2\pi\); la posición se
repite aunque el ángulo, que se registra de manera continua, no vuelva a cero.

En la circunferencia unidad,
\begin{equation}
  \mathbf r(t)
  =\bigl(\cos\omega t,\,\sin\omega t\bigr).
  \label{eq:circular-position}
\end{equation}
Al derivar, obtenemos
\begin{equation}
  \mathbf v(t)
  =\omega\bigl(-\sin\omega t,\,\cos\omega t\bigr).
  \label{eq:circular-velocity}
\end{equation}
La velocidad es tangente a la circunferencia. Podemos demostrar que es
perpendicular al radio sin basarnos en el dibujo:
\begin{align}
  \mathbf r\cdot\mathbf v
  &=\cos\omega t\,(-\omega\sin\omega t)
    +\sin\omega t\,(\omega\cos\omega t)\notag\\
  &=0.
  \label{eq:r-dot-v}
\end{align}

\begin{figure}[htbp]
  \centering
  \includegraphics[width=0.86\textwidth]{lecture_01_figure_12.png}
  \caption{Componentes de la posición y la velocidad para el movimiento sobre la
  circunferencia unidad. La cancelación en \(\mathbf r\cdot\mathbf v\) confirma que la velocidad
  es tangente a la circunferencia.}
  \lectureframe{01:24:30}
  \label{fig:circular-velocity-board}
\end{figure}

Al derivar una vez más,
\begin{equation}
  \mathbf a(t)
  =-\omega^2
  \bigl(\cos\omega t,\,\sin\omega t\bigr)
  =-\omega^2\mathbf r(t).
  \label{eq:centripetal-acceleration}
\end{equation}
La aceleración es paralela al radio, pero apunta en sentido contrario, hacia el
centro. Es la aceleración centrípeta.\footnote{Nota terminológica: la palabra
``centrífuga'' aparece en este punto de la clase, pero la ecuación y la
descripción adjunta identifican la aceleración centrípeta estándar, dirigida
hacia el centro.}

\begin{figure}[H]
  \centering
  \includegraphics[width=0.86\textwidth]{lecture_01_figure_13.png}
  \caption{Cálculo por componentes de la aceleración circular. Cada componente
  es \(-\omega^2\) por la componente correspondiente de la posición, así que
  \(\mathbf a=-\omega^2\mathbf r\).}
  \lectureframe{01:26:30}
  \label{fig:circular-acceleration-board}
\end{figure}

Como el radio es uno,
\begin{equation}
  \lvert\mathbf r\rvert=1,
  \qquad
  \lvert\mathbf v\rvert=\omega,
  \qquad
  \lvert\mathbf a\rvert=\omega^2.
\end{equation}
Aumentar \(\omega\) incrementa linealmente la rapidez y cuadráticamente la
aceleración dirigida hacia dentro que se requiere. Para una circunferencia de
radio \(R\), el mismo cálculo daría \(v=\omega R\) y
\(\lvert\mathbf a\rvert=\omega^2R=v^2/R\); en la clase se usa \(R=1\) para mantener
más sencilla la primera cuenta.

Ya disponemos del lenguaje elemental necesario para la mecánica de partículas:
estados y datos iniciales, evolución reversible, coordenadas, vectores y las
derivadas que definen la velocidad y la aceleración. La tarea siguiente es
determinar exactamente qué datos iniciales especifican una partícula y hallar
leyes que hagan evolucionar esos datos de un instante al siguiente.